% =====================================================================
%  Chapter 6 — The Axion
% =====================================================================

\chapter{The Axion and a Different Problem}
\label{ch:axion}

The strongest candidate in the field is not a candidate for dark matter. It
is a candidate for the solution to a different problem, which happens to
produce dark matter as a by-product. That distinction is the source of both
its strength and its difficulty.

\section{Why the axion was invented}The strong nuclear force appears to violate CP
symmetry, and there is no reason for it to do so. The Peccei--Quinn
mechanism resolves this by introducing a global symmetry, the Peccei--Quinn
symmetry, and a field, the axion, whose value relaxes the vacuum and
therefore removes the CP-violating term from the Lagrangian entirely.

The residual coupling of the axion to electromagnetism, the product
$g_{a\gamma\gamma}$, is predicted by the anomaly structure and is extremely
small, of order $10^{-16}$ times the fine structure constant in natural
units. For an axion with a mass below a few meV, that coupling is large
enough to be converted into a microwave signal --- a photon, at the axion
rest-frame frequency set by the axion mass, with a conversion probability
proportional to $g_{a\gamma\gamma}^2$.

This is the key structural fact about the axion: its small mass is what makes
it detectable, and its small mass is what the theory predicts for a solution
to the strong CP problem.

\section{The mass-coupling band}The relationship between the axion mass and its coupling
depends on the cosmological history, and the two standard scenarios differ
substantially.

In the \emph{pre-inflationary} scenario, the Peccei--Quinn symmetry is broken
before inflation and not restored afterwards, so the initial misalignment
angle is a random variable rather than a fixed one. The axion field behaves
like cold dark matter, the relic density is set by the misalignment mechanism,
and the relation between mass and coupling is a broad band with a
considerable width in the logarithm of the coupling.

In the \emph{post-inflationary} scenario, the symmetry is restored after
inflation and the misalignment angles are redistributed, giving a sharper
prediction: the relic density is set by the mean misalignment angle, and the
allowed band in the mass--coupling plane is narrow.

The practical consequence is that the axion occupies a band rather than a
point in parameter space, and the width of that band is a function of
cosmological history rather than of particle physics. Searches have to cover
the band, which is why the haloscope programme has been running for well over
a decade without reaching the QCD axion line in the canonical DFSZ scenario.

\section{Where the searches stand}ADMX has excluded the DFSZ axion at 3.3~$\mu$eV
\idx{ADMX Collaboration}, and has excluded the 1.1 to 1.3~GHz band. The
experiment works by converting axions into photons in a superconducting
cavity in a strong magnetic field, and the frequency of the expected signal
is directly the axion mass.

The haloscope programme is not the only approach, and the diversity is
necessary because the band is large.

\begin{itemize}
\item \textbf{Haloscones} --- dielectric haloscopes, and the development of
high-temperature superconductors, extend the reachable mass range downward.
\item \textbf{Helioscopes}, using the Sun as a magnetic-field source, search
for axion--photon conversion in the solar axion flux, which covers a
different and wider band than any terrestrial resonance experiment.
\item \textbf{Quantum sensors} --- single-photon detection and SQUID
devices are being developed for the same conversion physics with a
different noise floor.
\item \textbf{Astrophysical searches} --- axion conversion in magnetic fields
around neutron stars, and the axion's effect on stellar evolution, provide
constraints over a band no laboratory can reach.
\end{itemize}

\section{The case for and against}The case for the axion is strong. Its motivation is
independent of dark matter, its relic density is produced by a mechanism
that requires no fine tuning beyond an initial angle of order one, it is
stable against every decay channel in the Standard Model, and it is
technically within reach of experiments that exist.

The case against is equally specific. The coupling to photons depends on
whether the Peccei--Quinn symmetry is broken before or after
electroweak symmetry breaking, giving the DFSZ and KSVZ models, which
differ by roughly an order of magnitude in $g_{a\gamma\gamma}$ at a given
mass. Searches that are sensitive to one are not sensitive to the other.
And the dark matter interpretation requires either a post-inflationary
scenario, or a pre-inflationary one with a local misalignment angle that
happens to be favourable --- an assumption that cannot be checked from
inside the Solar System.

\section{Decision criteria}The axion would be confirmed by a photon excess at the mass
predicted for a specific cosmological scenario, at a significance that
survives the look-elsewhere correction over the band, seen by two
experiments. It would be excluded, in the region searched, by a null at the
same sensitivity. The QCD axion is not close to being excluded; the
canonical DFSZ band is the hardest target in the field and the most
valuable.

\begin{keytakeaways}
\item The axion is the strongest candidate because it solves the strong CP
problem first and produces dark matter as a by-product, not the other way
round.
\item The mass--coupling relation is a band whose width is set by
cosmological history, which means the search space is much larger than a
single experiment can cover.
\item DFSZ and KSVZ differ by an order of magnitude in coupling at a given
mass, and the two scenarios cannot both be covered by one instrument.
\item The canonical band is the hardest target in the field and the one with
the most to lose and the most to gain from a discovery.
\end{keytakeaways}
